% Generated draft: requires mathematical and source review.
\documentclass[11pt]{article}
\usepackage[margin=25mm]{geometry}
\usepackage{fontspec}
\setmainfont{Latin Modern Roman}
\usepackage{amsmath,amssymb,mathtools,unicode-math}
\setmathfont{Latin Modern Math}
\usepackage{enumitem,needspace,xurl,hyperref}
\setlength{\parindent}{0pt}
\setlength{\parskip}{5pt}
\setlength{\emergencystretch}{3em}
\newcommand{\N}{\mathbb{N}}
\newcommand{\Z}{\mathbb{Z}}
\newcommand{\Q}{\mathbb{Q}}
\newcommand{\R}{\mathbb{R}}
\newcommand{\C}{\mathbb{C}}
\newcommand{\F}{\mathbb{F}}
\newcommand{\A}{\mathbb{A}}
\newcommand{\PP}{\mathbb{P}}
\newcommand{\E}{\mathbb{E}}
\newcommand{\eps}{\varepsilon}
\newcommand{\dd}{\,\mathrm{d}}
\begin{document}
\section*{UnsolvedMath ID: 2708}
\textbf{Source problem code:} catalogue entry
\subsection{Definitions and mathematical statement}
A knot is an embedded circle $K\subset S^3$. A smooth concordance between knots is a properly embedded annulus $A\subset S^3\times[0,1]$ with the specified knots on its boundary. In the topological category the source requires local flatness of $A$ and local linearity of all group actions. Knots with involutions are pairs $(K,\tau)$ with $\tau^2=1$ and $\tau(K)=K$. A strongly invertible symmetry is orientation preserving on $S^3$ and reverses the knot, with a fixed circle meeting the knot in two points. A strongly negative amphichiral symmetry is orientation reversing on $S^3$, reverses the knot, and has two fixed points on it. Periodic symmetries have an axis disjoint from the knot; freely periodic symmetries act freely on $S^3$. Their finite cyclic orders must be fixed when an equivariant problem is specified.

For two knots with the same involution type, general equivariant concordance means an annulus $A$ invariant as a set under an involution $T$ of $S^3\times[0,1]$ extending the two endpoint involutions. Standard equivariant concordance additionally requires that $T$ preserve every level and restrict to a standard involution on each $S^3\times\{t\}$, as stated in the source. General $T$ need not preserve levels. Ordinary concordance requires no $T$. These definitions are considered separately in the smooth and locally flat, locally linear topological categories; no equivalence between them is assumed.

The five-part source problem is:
\begin{enumerate}
\item Describe the equivariant concordance groups, especially for strongly invertible and strongly negative amphichiral knots, using the appropriate based or directed symmetry data and equivariant connected sum. The source notes that the connected-sum point is one of the two fixed points on the knot and that dependence on this choice can obstruct commutativity. It does not give the complete auxiliary-data or category convention, so this structural request is not a single fully specified group-classification theorem.
\item Is the strongly negative amphichiral concordance group abelian? Once the group's conventions are fixed, this asks whether every pair of its classes commute under equivariant connected sum.
\item For a specified type of knot involution, can one find $(K_1,\tau_1)$ and $(K_2,\tau_2)$ that are generally equivariantly concordant but not standardly equivariantly concordant? The question can be asked for each type and in either category. Merely ordinarily concordant examples do not satisfy its stronger antecedent.
\item Is there a freely periodic knot that bounds an ordinary slice disk in $B^4$ but no invariant slice disk for an extension of its given cyclic symmetry? Equivariant disks/actions are smooth or locally flat/locally linear according to the chosen category. The source does not fix a period or explicitly distinguish standard from arbitrary ambient extensions for this slice question; that convention requires confirmation.
\item Boyle--Rouse's question asks whether every periodic or freely periodic $L$-space knot is a torus knot or an iterated torus knot. An $L$-space is a rational homology three-sphere $Y$ with $\dim_{\F_2}\widehat{HF}(Y;\F_2)=|H_1(Y;\Z)|$, and an $L$-space knot admits a nontrivial Dehn surgery yielding such a manifold. An iterated torus knot is obtained by successive cabling starting from a torus knot; no extra restrictions beyond the source's symmetry and $L$-space conditions are inserted.
\end{enumerate}

For (a), the source records Sakuma's homomorphism to the additive group of $\Z[t,t^{-1}]$, Di Prisa's theorem that the directed strongly invertible equivariant concordance group is nonabelian, and further invariants including a 2024 equivariant algebraic concordance theory whose precise reference is not supplied. Nonabelianity does not answer (b). Strongly negative amphichiral determinant-one examples of Boyle--Chen are motivated by the possibility that branched covers yield nontrivial $2$-torsion in the integer homology cobordism group; the source does not assert such nontrivial torsion is thereby established.

For (c), Sumners's involutions on $S^4$ with knotted fixed $S^2$ give, after puncturing, nonstandard involutions on $S^3\times[0,1]$ with standard fixed-circle endpoint actions, relevant to $2$-periodic and strongly invertible types. An ambient action alone does not supply the required invariant annulus counterexample. The source says interesting analogous nonstandard actions are not known there for freely periodic or strongly amphichiral types. For (d), Manolescu--Willis give freely $2$-periodic knots that are ordinarily concordant but not standardly equivariantly concordant. The source explicitly warns that they neither establish general equivariant concordance for (c) nor settle the freely periodic slice question.

For (e), Murasugi and Hartley give Alexander-polynomial factorization restrictions for periodic/free-periodic knots, while Ozsv\'ath--Szab\'o and Hedden--Watson give $L$-space restrictions. The background says the only known polynomials satisfying all these restrictions are products of cyclotomic polynomials; it does not promote this observation to a theorem about every such knot. Li--Ni's implication from a cyclotomic-product Alexander polynomial to torus/iterated-torus type is a conjecture. The newer Boyle--Rouse--Williams version proves the cyclotomic-product conclusion for freely periodic $L$-space knots of genus at most $16$, not the complete knot classification.

The literature triage checked 17 August 2026 classifies partial progress: meaningful strongly invertible structure is known, but (b), the stronger examples in (c), (d), and the general classification (e) remain open. The registry's main label is open. All these are supplied dated assessments.
\subsection{Short English statement}
Study equivariant knot concordance groups, their commutativity, general versus standard symmetric concordance, freely periodic sliceness, and symmetric $L$-space knots. Known nonabelianity and ordinary-concordance examples answer only parts of the agenda; the stronger symmetry and slice conditions must be preserved.
\subsection{Sources}
\begin{itemize}
\item Local UnsolvedMath registry, permanent ID $2708$, exact source code \texttt{catalogue entry}, ``Kirby Problem 1.49,'' including the supplied literature triage checked 17 August 2026.
\item R. \.{I}nan\c c Baykur, Robion C. Kirby and Daniel Ruberman, editors, \emph{K3: A New Problem List in Low-Dimensional Topology}, AMS Mathematical Surveys and Monographs \textbf{295} (2026), Problem 1.49, \url{https://math.berkeley.edu/sites/default/files/surv-295-ruberman-watermarked-author-pdf.pdf}.
\item [Sak86] Makoto Sakuma, \emph{On strongly invertible knots}, Algebraic and topological theories (Kinosaki, 1984), 176--196, Kinokuniya, Tokyo, 1986.
\item [DP23b] Alessio Di Prisa, \emph{The equivariant concordance group is not abelian}, Bulletin of the London Mathematical Society \textbf{55}(1) (2023), 502--507, \url{https://doi.org/10.1112/blms.12741}.
\item [BC24b] Keegan Boyle and Wenzhao Chen, \emph{Equivariant topological slice disks and negative amphichiral knots}, Indiana University Mathematics Journal \textbf{73}(5) (2024), 1623--1637.
\item [Sum75] D. W. Sumners, \emph{Smooth $\Z_p$-actions on spheres which leave knots pointwise fixed}, Transactions of the American Mathematical Society \textbf{205} (1975), 193--203, DOI 10.2307/1997199.
\item [BM23] Keegan Boyle and Jeffrey Musyt, \emph{Equivariant cobordisms between freely periodic knots}, Canadian Mathematical Bulletin \textbf{66}(2) (2023), 450--457, DOI 10.4153/S000843952200042X; part (d) is identified as its Question 1.
\item [MW25] Ciprian Manolescu and Michael Willis, \emph{A Rasmussen invariant for links in $\mathbb{RP}^3$}, Transactions of the American Mathematical Society, Series B \textbf{12} (2025), 789--830, DOI 10.1090/btran/221, \url{https://arxiv.org/abs/2301.09764}.
\item [Mur71] Kunio Murasugi, \emph{On periodic knots}, Commentarii Mathematici Helvetici \textbf{46} (1971), 162--174, DOI 10.1007/BF02566836.
\item [Har81] Richard Hartley, \emph{Knots with free period}, Canadian Journal of Mathematics \textbf{33}(1) (1981), 91--102, DOI 10.4153/CJM-1981-009-7.
\item [OS05b] Peter Ozsv\'ath and Zolt\'an Szab\'o, \emph{On knot Floer homology and lens space surgeries}, Topology \textbf{44}(6) (2005), 1281--1300, DOI 10.1016/j.top.2005.05.001.
\item [HW18] Matthew Hedden and Liam Watson, \emph{On the geography and botany of knot Floer homology}, Selecta Mathematica (N.S.) \textbf{24}(2) (2018), 997--1037, DOI 10.1007/s00029-017-0351-5.
\item [BR23b] Keegan Boyle and Nicholas Rouse, \emph{Obstructions to free periodicity and symmetric L-space knots}, 2023 preprint; the supplied triage identifies the updated version as Keegan Boyle, Nicholas Rouse and Ben Williams, arXiv:2310.01705v3 (2026), \url{https://arxiv.org/abs/2310.01705}.
\item [LN15] Eileen Li and Yi Ni, \emph{Half-integral finite surgeries on knots in $S^3$}, Annales de la Facult\'e des Sciences de Toulouse, Math\'ematiques (6) \textbf{24}(5) (2015), 1157--1178, DOI 10.5802/afst.1479.
\end{itemize}
\end{document}
